\documentclass[../../../include/open-logic-section]{subfiles}

\begin{document}

\olfileid{sth}{ord-arithmetic}{expo}
\olsection{ક્રમવાચક સંખ્યાઓની ઘાતક્રિયા}

હવે ક્રમવાચક સંખ્યાઓની ઘાતક્રિયા જોઈએ. દુર્ભાગ્યે,
તે માટે કોઈ \emph{સરસ} રચનાત્મક વ્યાખ્યા નથી.

અલબત્ત, સ્પષ્ટ રચનાત્મક વ્યાખ્યાઓ \emph{છે}.
એક આ છે. $\text{finfun}(\alpha,\beta)$ એ બધાં
એવાં વિધેયો $f
\colon \alpha \to \beta$નો ગણ હોય કે
$\Setabs{\gamma \in
\alpha}{f(\gamma) \neq 0}$ કોઈ પ્રાકૃતિક સંખ્યા
જેટલા ઘટકો ધરાવે. $\text{finfun}(\alpha,\beta)$
પર $f
\sqsubset g$ ત્યારે અને માત્ર ત્યારે એવો સુક્રમ
વ્યાખ્યાયિત કરો કે $f \neq g$ હોય અને
$f(\gamma_0) < g(\gamma_0)$, જ્યાં
$\gamma_0 = \text{max}\Setabs{\gamma \in \alpha}{f(\gamma) \neq
g(\gamma)}$. પછી $\ordexpo{\alpha}{\beta}$ને
$\ordtype{\text{finfun}(\beta, \alpha), \sqsubset}$
તરીકે વ્યાખ્યાયિત કરી શકીએ. પોટર આ સ્પષ્ટ વ્યાખ્યા
વાપરે છે અને તરત જ સમજાવે છે:
\sourcecorrection{OLMD-132}{મૂળમાં
$\text{finfun}(\alpha,\beta)$ એ $\alpha$થી
$\beta$ સુધીનાં મર્યાદિત આધારવાળાં વિધેયો છે,
પરંતુ તેનો ક્રમપ્રકાર $\ordexpo{\alpha}{\beta}$
લખાયો છે. આ વિધેયોનો ઊલટો-શબ્દકોશક્રમ
$\ordexpo{\beta}{\alpha}$ આપે છે. તેથી અહીં
$\ordexpo{\alpha}{\beta}$ માટે
$\text{finfun}(\beta,\alpha)$ લીધો છે; છેલ્લાં
સ્વાધ્યાયમાં પણ એ જ સુધારો છે.}
\begin{quote}
	આ ક્રમની પસંદગી માત્ર યોગ્ય પુનરાવર્તી શરત
	સંતોષતી ક્રમવાચક ઘાતક્રિયાની વ્યાખ્યા મેળવવાની
	આપણી ઇચ્છાથી નક્કી થાય છે\ldots; અને ક્રમિત
	સરવાળા કે ક્રમિત ગુણાકાર કરતાં તેનું ચિત્ર મનમાં
	ઊભું કરવું ઘણું અઘરું છે.
	\citep[p.~199]{Potter2004}
\end{quote}
બરાબર. સરવાળાને ``$\alpha$ની નકલ પછી
$\beta$ની નકલ'' તરીકે અને ગુણાકારને
``$\alpha$ની નકલોનો $\beta$-અનુક્રમ'' તરીકે
સમજાવ્યા હતા. પરંતુ $\text{finfun}(\alpha, \gamma)$
વિશે એવું કોઈ ટૂંકું સહજ વર્ણન નથી. તેથી ક્રમવાચક
ઘાતક્રિયાની વ્યાખ્યા \emph{સીધી} અનંતાતીત
પુનરાવર્તનથી આપીએ છીએ, એટલે કે:

\begin{defn}\ollabel{ordexporecursion}
\begin{align*}
	\ordexpo{\alpha}{0} &= 1\\
	\ordexpo{\alpha}{\beta\ordplus 1} &=\ordexpo{\alpha}{\beta} \ordtimes \alpha\\
	\ordexpo{\alpha}{\beta} &= \bigcup_{\delta < \beta}\ordexpo{\alpha}{\delta}& & \text{when $\beta$ is a limit ordinal}
\end{align*}
\sourcecorrection{OLMD-133}{મૂળનું સીમા-સમીકરણ
આધાર $\alpha=0$ માટે ઉપરની રચનાત્મક વ્યાખ્યાથી
મેળ ખાતું નથી. $\ordexpo{0}{0}=1$ અને
પછીના અનુગામી ઘાત $0$ હોવાથી પ્રદર્શિત સીમા-યોગગણ
$\ordexpo{0}{\omega}=1$ આપે છે; પરંતુ
$\text{finfun}(\omega,0)$ ખાલી હોવાથી તેનો
ક્રમપ્રકાર $0$ છે. શૂન્ય આધાર માટે અલગ શરત વિના
આ પ્રદર્શનને બંને વ્યાખ્યાઓની સમકક્ષતા ગણવી નહીં.}
\end{defn}

આપણે \emph{ગણસિદ્ધાંતીઓ તરીકે} કામ કરતા હોત,
તો ક્રમવાચક ઘાતક્રિયાના કેટલાક ગુણધર્મો વધુ
તપાસતા. પરંતુ તે અદલબદલીય નથી એ અપેક્ષિત
તથ્ય સિવાય અહીં વધુ ઉમેરવાનું નથી. ખરેખર
$\ordexpo{2}{\omega} =
\bigcup_{\delta < \omega}\ordexpo{2}{\delta} = \omega$,
જ્યારે $\ordexpo{\omega}{2} = \omega \ordtimes \omega$.
પણ ઘાતક્રિયા અદલબદલીય હશે એવી \emph{અપેક્ષા}
રાખવી જ ન જોઈએ, કારણ કે પ્રાકૃતિક સંખ્યાઓમાં પણ
એવું નથી: $\ordexpo{2}{3} = 8 < 9 = \ordexpo{3}{2}$.

\begin{prob}
અનંતાતીત અનુમાનથી સાબિત કરો કે જો
$\ordexpo{\alpha}{\beta} = \ordtype{\text{finfun}(\beta, \alpha),
\sqsubset}$ વ્યાખ્યાયિત કરીએ, તો
\olref[sth][ord-arithmetic][expo]{ordexporecursion}નાં
પુનરાવર્તી સમીકરણો મળે છે.
\sourcecorrection{OLMD-134}{મૂળના આ સ્વાધ્યાયમાં
પણ OLMD-132 મુજબ વિધેયોના દલીલોનું સ્થાન
સુધાર્યું છે. તેમ છતાં OLMD-133 મુજબ $\alpha=0$
અને સીમા $\beta$ માટે પ્રદર્શિત સમીકરણ ખોટું પડે છે;
તેથી એ શરત સુધાર્યા વિના આ સ્વાધ્યાય બધા
ક્રમવાચક આધાર માટે સાબિત કરી શકાતો નથી.}
\end{prob}

\end{document}
